\documentclass[a4paper,11pt,center]{article}
\usepackage[utf8]{inputenc}
\usepackage{graphicx}

\begin{document}
	
	
	\begin{center}
		\textbf{MOUNTAINS OF THE MOON UNIVERSITY\newline\newline\newline
			FACULTY OF SCIENCE TECHNOLOGY\newline\newline\newline
			AND INNOVATION\newline\newline\newline
			DEPARTMENT OF COMPUTER SCIENCE\newline\newline\newline
			Individual Coursework}\newline\newline\newline
	\end{center}
	
	
	
	\textbf{NAME: Ushindi Edward}\newline\newline
	\textbf{REGISTRATION NUMBER: 230270/U/MMU/BCS/00530}\newline\newline
	\textbf{COURSE CODE: BCS 1201}\newline\newline
	\textbf{LECTURER'S NAME: Mr. Samuel Ocen}
	
	
	
	
	
	
	\newpage
	
	
	\section{Resource Allocation in Cloud 
		Computing}
	\subsection{\textbf{Number One: Basic Resource Allocation}}

	
	\begin{verbatim}
		# NUMBER 1
		
		
		# import necessary libraries
		from pulp import LpVariable, LpProblem,LpMinimize
		
		# define the linear problem
		lp = LpProblem(name="minimizing_cost", sense=LpMinimize)
		
		# define the decision variables
		x = LpVariable(name="X")
		y = LpVariable(name="Y")
		
		#define the objectives
		lp += 4*x + 5*y ,"objective"
		
		# define constraints
		lp += 2*x + 3*y>=10, "CPU"
		lp += x + 2*y>=5, "Memory"
		lp += 3*x + y>=8, "Storage"
		
		# solve the results
		lp.solve()
		
		# print results
		print("OPTIMUM SOLUTION")
		print(f"X = {x.varValue}")
		print(f"Y = {y.varValue}")
		print(f"optimum solution: {lp.objective.value()}")
		
		# the graph
		#import libraries
		import numpy as np
		from scipy.optimize import linprog
		import matplotlib.pyplot as plt
		#x array
		x=np.linspace(0,10,100)
		#convert constraints to inequalities
		y1 = (10-2*x)/3
		y2 = (5-x)/2
		y3 = (8-3*x)
		#plot constraits
		plt.plot(x,y1, label="2x + 3y >=10 (cpu)")
		plt.plot(x,y2, label="x + 2y >=5 (memory)")
		plt.plot(x,y3, label="3x + y >=8 (storage)")
		#feasible region
		plt.fill_between(x,0,np.minimum.reduce([y1,y2,y3]),color="blue",alpha=0.5,label="feasible region")
		plt.xlabel("x")
		plt.ylabel("y")
		plt.title("feasible region for Basic Resource Allocation")
		plt.ylim(0,5)
		plt.xlim(0,5)
		plt.legend()
		plt.show()
	\end{verbatim}
	
	\textbf{Solution}\newline\newline
	
	\includegraphics{C:/Users/skurge/OneDrive/Desktop/USHINDI EDWARD BCS 00530/num1.png}\newline\newline
	
	\newpage
	\textbf{Graph}\newline\newline
	\includegraphics{graph1}
	

	\newpage
	
	
	
	
	\subsection{\textbf{Number Two: Load Balancing}}

	
	\begin{verbatim}
		# NUMBER 2
		
		
		# import libraries
		from pulp import *
		
		# define lp
		lp = LpProblem(name="minimizing_the_overall_response_time", sense=LpMinimize)
		
		# define decision variables
		x = LpVariable(name="x", lowBound=0)
		y = LpVariable(name="y", lowBound=0)
		
		# define the objective
		lp += 5*x + 4*y, "objective"
		
		# define constraints
		lp += 2*x + 3*y <=20, "server1"
		lp += 4*x + 2*y <=15, "server2"
		
		# solve
		lp.solve()
		
		print("RESULTS")
		print(f"x = {x.varValue}")
		print(f"y = {y.varValue}")
		print(f"optimum solution = {lp.objective.value()}")
		
		
		# the graph
		#import libraries
		import numpy as np
		from scipy.optimize import linprog
		import matplotlib.pyplot as plt
		#x array
		x=np.linspace(0,10,100)
		#convert constraints to inequalities
		y1 = (20-2*x)/3
		y2 = (15-4*x)/2
		#plot constraits
		plt.plot(x,y1, label="2x + 3y <=20 (server1)")
		plt.plot(x,y2, label="4x + 2y <=15 (server2)")
		#feasible region
		plt.fill_between(x,0,np.minimum(y1,y2),color="grey",alpha=0.5,label="feasible region")
		plt.xlabel("x")
		plt.ylabel("y")
		plt.title("feasible region for Load balancing")
		plt.ylim(0,10)
		plt.xlim(0,10)
		plt.legend()
		plt.show()		
	\end{verbatim}
	
	\textbf{Solution}\newline\newline
	\includegraphics{num2}\newline\newline
	\includegraphics{graph2}

	\newpage
	
	
	
	
	\subsection{\textbf{Number Three: Energy Efficient Resource Allocation}}

	\begin{verbatim}
		# NUMBER 3
		
		
		# import necessary libraries
		from pulp import *
		
		# define the linear problem
		Lp = LpProblem(name="cloud data", sense=LpMinimize)
		
		#  define the decision values
		x = LpVariable(name="x", lowBound=0)
		y = LpVariable(name="y", lowBound=0)
		
		# define the objectives
		Lp += 3*x + 2*y, "objective"
		
		# define the constraints
		Lp += 2*x + 3*y >=15, "cpu_allocation"
		Lp += 4*x + 2*y >=10, "memory_allocation"
		
		# solve the results
		Lp.solve()
		
		# print results
		print("Result")
		print(f"X = {x.varValue}")
		print(f"Y = {y.varValue}")
		print(f"optimum solution: {Lp.objective.value()}")
		
		# the graph
		#import libraries
		import numpy as np
		from scipy.optimize import linprog
		import matplotlib.pyplot as plt
		#x array
		x=np.linspace(0,20,100)
		#convert constraints to inequalities
		y1 = (15-2*x)/3
		y2 = (10-4*x)/2
		#plot constraits
		plt.plot(x,y1, label="2x + 3y >=15 (cpu allocation)")
		plt.plot(x,y2, label="4x + 2y >=10 (memory allocation")
		#feasible region
		plt.fill_between(x,100,np.maximum.reduce([y1,y2]),color="orange",alpha=0.5,label="feasible region")
		plt.xlabel("x")
		plt.ylabel("y")
		plt.title("feasible region for energy Efficient resource allocation")
		plt.ylim(0,15)
		plt.xlim(0,10)
		plt.legend()
		plt.show()
		
	\end{verbatim}
	
	\textbf{Solution}\newline\newline
	\includegraphics{num3}\newline\newline
	\includegraphics{graph3}
	\newpage
	
	
	
	
	\subsection{\textbf{Number Four: Multi-Tenant Resource Sharing}}

	\begin{verbatim}
		# NUMBER 4
		
		
		# import libraries
		from pulp import *
		
		# define lp
		lp = LpProblem(name="allocating_resources", sense=LpMinimize)
		
		# define decision variables
		x = LpVariable(name="x", lowBound=0)
		y = LpVariable(name="y", lowBound=0)
		
		# define the objective
		lp += 5*x + 4*y, "objective"
		
		# define constraints
		lp += 2*x + 3*y >=12, "tenant 1"
		lp += 4*x + 2*y >=18, "tenant 2"
		
		# solve
		lp.solve()
		
		print("RESULTS")
		print(f"x = {x.varValue}")
		print(f"y = {y.varValue}")
		print(f"optimum solution = {lp.objective.value()}")
		
		# the graph
		#import libraries
		import numpy as np
		from scipy.optimize import linprog
		import matplotlib.pyplot as plt
		#x array
		x=np.linspace(0,25,100)
		#convert constraints to inequalities
		y1 = (12-2*x)/3
		y2 = (18-4*x)/2
		#plot constraits
		plt.plot(x,y1, label="2x + 3y >=12 (tenant1)")
		plt.plot(x,y2, label="4x + 2y >=18 (tenant2)")
		#feasible region
		plt.fill_between(x,200,np.maximum.reduce([y1,y2]),color="red",alpha=0.5,label="feasible region")
		plt.xlabel("x")
		plt.ylabel("y")
		plt.title("feasible region for Multi_Tenant Resource Sharing")
		plt.ylim(0,20)
		plt.xlim(0,20)
		plt.legend()
		plt.show()
		
	\end{verbatim}
	
	\textbf{Solution}\newline\newline
	\includegraphics{num4}\newline\newline
	\includegraphics{graph4}
	
	\newpage
	
	
	
	
	\section{Transportation Logistics 
		Optimization}
	\subsection{\textbf{Number Five: Production Planning in Manufacturing}}

	\begin{verbatim}
		# NUMBER 5
		
		
		# import libraries
		from pulp import *
		
		# define lp
		lp = LpProblem(name="production_costs", sense=LpMinimize)
		
		# define decision variables
		x1 = LpVariable(name="x1", lowBound=0)
		x2= LpVariable(name="x2", lowBound=0)
		x3= LpVariable(name="x3", lowBound=0)
		
		# define the objective
		lp += 5*x1 + 3*x2 + 4*x3, "objective"
		
		# define constraints
		lp += 2*x1 + 3*x2 + 1*x3 <=1000, "raw material"
		lp += 4*x1 + 2*x2 + 5*x3 <=120, "labor hours"
		lp += x1 >=200
		lp += x2 >=300
		lp += x3 >=150
		
		# solve
		lp.solve()
		
		print("RESULTS")
		print(f"x1 = {x1.varValue}")
		print(f"x2 = {x2.varValue}")
		print(f"x3 = {x3.varValue}")
		print(f"optimum solution x1,x2,x3 = {lp.objective.value()}")
		
		
		# Plotting the feasible region (in 3D space)
		fig = plt.figure()
		ax = fig.add_subplot(111, projection='3d')
		
		x = np.linspace(0, 400, 100)
		y = np.linspace(0, 400, 100)
		x, y = np.meshgrid(x, y)
		z1 = (1000 - 2*x - 3*y) / 1
		z2 = (120 - 4*x - 2*y) / 5
		
		ax.plot_surface(x, y, np.maximum(z1, z2), alpha=0.5, rstride=100, cstride=100, color='gray', label='Feasible Region')
		ax.set_xlabel('x1')
		ax.set_ylabel('x2')
		ax.set_zlabel('x3')
		ax.set_title('Feasible Region for Production Planning in manufacturing')
		ax.scatter(x[0], x[1], x[2], color='orange', label='Optimal Solution')
		
		plt.show()
	\end{verbatim}

	\textbf{Solution}\newline\newline
	\includegraphics{num5}\newline\newline
	\includegraphics{graph5}

	
	\newpage
	
	
	
	
	\subsection{\textbf{Number Six: Financial Portfolio Optimization}}

	
	\begin{verbatim}
		#libraries
		import numpy as np
		import matplotlib.pyplot as plt
		from mpl_toolkits.mplot3d import Axes3D
		from pulp import LpVariable, LpMaximize, LpProblem
		
		#the problem
		problem = LpProblem(name="Financial_portfolio_optimization", sense=LpMaximize)
		
		#decision variables
		x1 = LpVariable(name="x", lowBound=0)
		x2 = LpVariable(name="y", lowBound=0)
		x3 = LpVariable(name="z", lowBound=0)
		
		# Define the objective function coefficients
		problem += 0.08*x1 + 0.1*x2 + 0.12*x3, "objective"
		
		# Coefficients of the inequality constraints
		problem += 2*x1 + 3*x2 + x3 <= 10000, "budget"
		problem += x1 >= 2000
		problem += x2 >= 1500
		problem += x3 >= 1000
		
		
		# Solve linear programming problem
		problem.solve()
		
		# Display the results
		print("OPTIMAL SOLUTION:")
		print(f"X: {x1.varValue}")
		print(f"Y: {x2.varValue}")
		print(f"Z: {x3.varValue}")
		print(f"Minimum cost: {problem.objective.value()}")
		
		
		# Create a meshgrid for A, B, and C
		A_vals = np.linspace(0, 2500, 50)
		B_vals = np.linspace(0, 2000, 50)
		A_grid, B_grid = np.meshgrid(A_vals, B_vals)
		
		# Calculate the corresponding z-values (ROI function)
		C_vals = (10000 - 2 * A_grid - 3 * B_grid)  # Budget constraint
		ROI_vals = 0.08 * A_grid + 0.1 * B_grid + 0.12 * C_vals
		
		# Create the 3D plot
		fig = plt.figure(figsize=(10, 6))
		ax = fig.add_subplot(111, projection='3d')
		
		# Plot the feasible region (constraints)
		ax.plot([2000, 2000], [0, 2000], [0, 470], color='red', linestyle='--', linewidth=2, label='Budget')
		ax.plot([0, 2500], [1500, 1500], [0, 470], color='green', linestyle='--', linewidth=2, label='Maximum Investment')
		ax.plot([0, 2500], [0, 2000], [470, 470], color='blue', linestyle='--', linewidth=2, label='Maximum Investment')
		
		# Highlight the optimum point
		optimum_A = 2000
		optimum_B = 1500
		optimum_ROI = 470
		ax.the plotscatter(optimum_A, optimum_B, optimum_ROI, color='purple', s=100, label='Optimum Point')
		
		# Set labels and title
		ax.set_xlabel('A')
		ax.set_ylabel('B')
		ax.set_zlabel('ROI')
		ax.set_title('Return On Investiment Maximization')
		
		# Add a legend
		ax.legend()
		
		# Show 
		plt.show()
\end{verbatim}
		
	\newpage
	\textbf{Solution}\newline\newline
	\includegraphics{num6}\newline\newline
	\includegraphics{graph6}
	
	
	

	
	
	\newpage
	
	
	
	
	\section{\textbf{Project Resource Allocation 
			Optimization}}
	\subsection{\textbf{Number Seven: Diet Optimization}}

	
	\begin{verbatim}
		# NUMBER 7
		
		
		
		# import libraries
		from pulp import *
		
		# define lp
		lp = LpProblem(name="daily_diet", sense=LpMinimize)
		
		# define decision variables
		x1 = LpVariable(name="x1", lowBound=0)
		x2 = LpVariable(name="x2", lowBound=0)
		
		# define the objective
		lp += 3*x1 + 4*x2, "objective"
		
		# define constraints
		lp += 2*x1 + 1*x2 >=20, "protein"
		lp += 3*x1 + 2*x2 >=25, "vitamin"
		
		# solve
		lp.solve()
		
		print("RESULTS")
		print(f"x1 = {x1.varValue}")
		print(f"x2 = {x2.varValue}")
		print(f"optimum solution = {lp.objective.value()}")
		
		
		#graph
		#import libraries
		import matplotlib.pyplot as plt
		import numpy as np
		
		# Plotting the feasible region
		x1 = np.linspace(0, 20, 100)
		
		x2_1 = 20 - 2*x1  # Protein constraint
		x2_2 = (25 - 3*x1) / 2  # Vitamins constraint
		
		plt.plot(x1, x2_1, label='2x1 + x2 >= 20 (Protein)')
		plt.plot(x1, x2_2, label='3x1 + 2x2 >= 25 (Vitamins)')
		
		plt.fill_between(x1,200, np.maximum.reduce([x2_1,x2_2 ]), color='yellow', alpha=0.4, label='Feasible Region')
		plt.xlim(0,20)
		plt.ylim(0,30)
		plt.xlabel('x1 (Food Item 1)')
		plt.ylabel('x2 (Food Item 2)')
		plt.title('Feasible Region for Diet Optimization')
		plt.legend()
		plt.show()
	\end{verbatim}
	
	\textbf{Solution}\newline\newline
	\includegraphics{num7}\newline\newline
	\includegraphics{graph7}	
			
	\newpage
	
	
	
	
	\subsection{\textbf{Number Eight: Production Planning}}

	
	\begin{verbatim}
		# NUMBER 8
		
		
		
		# import libraries
		from pulp import *
		
		# define lp
		lp = LpProblem(name="profi_maximization", sense=LpMaximize)
		
		# define decision variables
		x1 = LpVariable(name="x1", lowBound=0)
		x2 = LpVariable(name="x2", lowBound=0)
		
		# define the objective
		lp += 5*x1 + 3*x2, "objective"
		
		# define constraints
		lp += 2*x1 + 3*x2 <=60, "labor"
		lp += 4*x1 + 2*x2 <=80, "raw materials"
		
		# solve
		lp.solve()
		
		print("RESULTS")
		print(f"x1 = {x1.varValue}")
		print(f"x2 = {x2.varValue}")
		print(f"optimum solution = {lp.objective.value()}")
		
		
		#graph
		#import libraries
		import matplotlib.pyplot as plt
		import numpy as np
		
		# Plotting the feasible region
		x1 = np.linspace(0, 100, 100)
		# making one of the variables the subject in each constrait
		x2_1 = (60 - 2*x1)/3  # labor
		x2_2 = (80 - 4*x1) / 2  # raw materials
		
		plt.plot(x1, x2_1, label='2x1 + 3x2 <= 60 (labor)')
		plt.plot(x1, x2_2, label='4x1 + 2x2 <= 80 ( raw materials)')
		
		plt.fill_between(x1,0, np.minimum.reduce([x2_1,x2_2] ), color='pink', alpha=0.6, label='Feasible Region')
		plt.xlim(0,50)
		plt.ylim(0,50)
		plt.xlabel('x1 (Labor)')
		plt.ylabel('x2 (raw materials)')
		plt.title('Feasible Region for Production planning')
		plt.legend()
		plt.show()

	\end{verbatim}
	
	\textbf{Solution}\newline\newline
	\includegraphics{num8}\newline\newline
	\includegraphics{graph8}

	
\end{document}